\section{The quantum model, uniqueness, and positive frames}
\label{s:quantum-frames}

\subsection{Polynomial normalization and the edge kernel}

We use the Pauli matrices
\[
 \sigma_1=\begin{pmatrix}0&1\\1&0\end{pmatrix},\qquad
 \sigma_2=\begin{pmatrix}0&-i\\i&0\end{pmatrix},\qquad
 \sigma_3=\begin{pmatrix}1&0\\0&-1\end{pmatrix},
 \qquad a\cdot\sigma=\sum_{j=1}^3a_j\sigma_j.
\]

\begin{definition}\label{d:poly}
Let \(V_d\) be the homogeneous binary polynomials of degree \(d\).
Its monomials \(z_1^{d-j}z_2^j\), \(0\le j\le d\), are orthogonal with
squared norms \(\binom dj^{-1}\). The action is
\((U_gp)(z)=p(g^{-1}z)\). Use \(V_4\) at each physical site.
For each of the finitely many frame letters \(a\) choose one unit
spinor \(\zeta(a)\in\mathbb C^2\) with
\(\zeta(a)\zeta(a)^*=(I+a\cdot\sigma)/2\).
No spinor choice on the whole sphere is part of the definitions.
\end{definition}

Here \(d\) is a nonnegative integer and \(g\in\mathrm{SU}(2)\).
We record explicitly why this is the spin-\(d/2\) Hilbert space.
In the \(d\)-fold tensor product of the dual of \(\mathbb C^2\),
average the \(\binom dj\) distinct elementary tensors having \(j\)
second-coordinate covectors. Evaluation of this average on
\(z^{\otimes d}\) is \(z_1^{d-j}z_2^j\), and its squared norm is
\(\binom dj\binom dj^{-2}=\binom dj^{-1}\).
These averages are orthogonal for different \(j\).
The dual action of a unitary matrix is unitary, as is its tensor
power; restriction to symmetric tensors therefore gives precisely
the invariant inner product of Definition~\ref{d:poly}.

For completeness, irreducibility follows directly from the monomials
\(m_j=z_1^{d-j}z_2^j\). The diagonal subgroup
\(g_t=\operatorname{diag}(e^{it/2},e^{-it/2})\) acts on \(m_j\)
by \(e^{it(j-d/2)}\). When \(d\ge1\), take
\(t=2\pi/(d+1)\); the \(d+1\) eigenvalues are distinct, so each
monomial projection is a polynomial in \(U_{g_t}\).
A nonzero invariant subspace consequently contains some \(m_j\).
For the skew-Hermitian trace-zero matrices
\[
 A=\begin{pmatrix}0&1\\-1&0\end{pmatrix},\qquad
 C=\begin{pmatrix}0&i\\i&0\end{pmatrix},
\]
differentiating the action along \(\exp(tA)\) and \(\exp(tC)\)
gives, respectively,
\(dU(A)=z_1\partial_{z_2}-z_2\partial_{z_1}\) and
\(dU(C)=-i(z_1\partial_{z_2}+z_2\partial_{z_1})\).
The subspace is stable under these operators and hence under their
complex linear combinations
\(z_1\partial_{z_2}=(dU(A)+i\,dU(C))/2\) and
\(z_2\partial_{z_1}=(-dU(A)+i\,dU(C))/2\).
Their actions are
\[
 z_1\partial_{z_2}m_j=j m_{j-1}\quad(j\ge1),\qquad
 z_2\partial_{z_1}m_j=(d-j)m_{j+1}\quad(j<d).
\]
The indicated coefficients are nonzero. Iterating reaches all
monomials, proving irreducibility. For \(d=0\) it is immediate.
The weights are \(-d/2,-d/2+1,\ldots,d/2\), which identifies the
spin. Differentiation here preserves the invariant subspace because
it is finite dimensional and hence closed.

If \(|a|=1\), the Hermitian matrix \((I+a\cdot\sigma)/2\) is a
rank-one orthogonal projection: its trace is one and
\((a\cdot\sigma)^2=I\). A unit vector in its range supplies the
spinor required in Definition~\ref{d:poly}. We make this choice only
for the finite sets under consideration. Changing such a choice by
a phase does not change any squared modulus below.
The tensor product \(\mathcal H_\Lambda\) is identified with
polynomials homogeneous of degree four in each pair
\(z_x=(z_{x,1},z_{x,2})\).

\begin{lemma}\label{l:edge}
For \(e=\{x,y\}\) and \(\varphi\in\mathcal H_\Lambda\),
\(P_e\varphi=0\) if and only if
\(\varphi=q_e\chi\), where \(\chi\) is homogeneous of degree three
in \(z_x,z_y\) and degree four in every other site variable.
\end{lemma}

\begin{proof}
First consider just the two endpoint sites. Define the diagonal
evaluation map
\[
 \Delta:V_4\otimes V_4\longrightarrow V_8,\qquad
 (\Delta p)(z)=p(z,z).
\]
It is equivariant for the diagonal \(\mathrm{SU}(2)\) action.
For \(0\le k\le8\), choose \(j=\max\{0,k-4\}\); then both
\(j\) and \(k-j\) lie in \(\{0,\ldots,4\}\), and
\[
 \Delta\bigl(z_{x,1}^{4-j}z_{x,2}^j
              z_{y,1}^{4-(k-j)}z_{y,2}^{k-j}\bigr)
       =z_1^{8-k}z_2^k.
\]
Thus \(\Delta\) is onto and
\(\dim\ker\Delta=25-9=16\).

Multiplication by \(q_{xy}\) injects \(V_3\otimes V_3\) into
\(V_4\otimes V_4\), since the polynomial ring is an integral
domain and \(q_{xy}\ne0\). Its image is contained in
\(\ker\Delta\), because \(q_{xy}(z,z)=0\). Its dimension is
\(4\cdot4=16\), so
\[
 \ker\Delta=q_{xy}(V_3\otimes V_3).
\]

We identify the orthogonal complement with the physical total-spin
space, using the inner products just fixed. Since \(\Delta\) is
onto, \(\Delta^*\) is injective. Unitarity and equivariance imply
that \(\Delta^*\) also intertwines the actions. Hence
\(\operatorname{Ran}\Delta^*=(\ker\Delta)^\perp\) is an
irreducible copy of \(V_8\), namely spin four.
In \(V_4\otimes V_4\), the weight of \(m_j\otimes m_k\) is
\(j+k-4\). Its weight-four subspace is one dimensional, spanned
by \(m_4\otimes m_4\). Every spin-four copy contains a nonzero
weight-four vector. Two such copies therefore intersect nontrivially,
and irreducibility forces them to coincide. Consequently
\(\operatorname{Ran}\Delta^*\) is the entire total-spin-four
space onto which \(P_e\) projects. Thus \(\ker P_e=\ker\Delta\).
Reversing the edge orientation changes the determinant by \(-1\)
and leaves its image unchanged.

Finally, expand a general \(\varphi\) in a tensor monomial basis
of the other sites. The equation \(P_e\varphi=0\) is equivalent
to the same equation for each endpoint coefficient. Each coefficient
is \(q_e\) times a polynomial of bidegree \((3,3)\), by the
two-site argument. Summing those quotients against the other-site
monomials gives exactly the stated \(\chi\). The converse follows
coefficient by coefficient as well.
\end{proof}

\begin{lemma}\label{l:divisibility}
Let \(L,N\ge3\), \(R\subset\Lambda_{L,N}\), and
\(E_R=\{e\in E_{L,N}:e\cap R\ne\varnothing\}\).
If \(\varphi\in\mathcal H_\Lambda\) and \(P_e\varphi=0\)
for every \(e\in E_R\), then
\(\varphi=(\prod_{e\in E_R}q_e)\chi\), where \(\chi\) is independent
of every \(z_x\), \(x\in R\).
\end{lemma}

\begin{proof}
The four sites \(x+(1,0),x-(1,0),x+(0,1),x-(0,1)\) are
different from \(x\) and pairwise distinct. Indeed, within a
coordinate the equality of the two neighbors would make its period
divide two; between coordinates an equality would make a period
divide one. Neither is possible for periods at least three.
Thus there are exactly four distinct incident edges at every site.

Work in the unique factorization domain
\(\mathcal R=\mathbb C[z_{u,1},z_{u,2}:u\in\Lambda]\).
For two distinct sites \(x,y\), regard
\(q_{xy}=z_{x,1}z_{y,2}-z_{x,2}z_{y,1}\) as a polynomial in
\(z_{y,2}\) over the polynomial ring in all the other variables.
Its two coefficients are relatively prime: the variable
\(z_{x,1}\) does not divide \(z_{x,2}z_{y,1}\).
It is primitive and is linear and nonconstant over the fraction
field of that coefficient ring. Gauss's lemma therefore makes it
irreducible in \(\mathcal R\). This argument already includes all
other site variables. Each \(q_e\), with either orientation, is
accordingly an irreducible and hence a prime element.

Two different unoriented edges have different endpoint sets.
Their determinants depend nontrivially on different sets of site
variables and cannot differ by a nonzero scalar, the only kind of
unit in \(\mathcal R\). They are therefore nonassociate.
Lemma~\ref{l:edge} says that every \(q_e\), \(e\in E_R\),
divides \(\varphi\). Unique factorization for these distinct prime
factors implies that their product, denoted \(b_R\), divides
\(\varphi\).

If \(\varphi=0\), take \(\chi=0\). Otherwise write
\(\varphi=b_R\chi\) and decompose \(\chi\) into its finitely
many multihomogeneous components. Write \(d_R(u)\) for the number
of edges of \(E_R\) at \(u\), so that \(b_R\) has degree
\(d_R(u)\) in \(z_u\). Multiplication by the nonzero \(b_R\)
is injective, and distinct multidegrees remain distinct after adding
\((d_R(u))_u\). Since \(\varphi\) has multidegree \((4)_u\),
the only possible component of \(\chi\) has degree
\(4-d_R(u)\) at each \(u\). For \(u\in R\), all four incident
edges are in \(E_R\), so this degree is zero. A polynomial of
degree zero in both variables at each such site is independent of
those variables. This also covers \(R=\varnothing\), where
\(b_R=1\) and the independence assertion is vacuous.
\end{proof}

\begin{proposition}\label{p:toruskernel}
For every \(L,N\ge3\), the torus is simple and has four distinct
neighbors at each site, \(B_{L,N}\ne0\), and
\(\ker H_{L,N}=\mathbb C B_{L,N}\).
Changing edge orientations does not change this one-dimensional space.
\end{proposition}

\begin{proof}
The neighbor calculation in Lemma~\ref{l:divisibility} proves
the asserted degree and the absence of loops, including at period
three; the edge set in Definition~\ref{d:model} has no repeated
unordered pairs. Each determinant is nonzero, so their product is
nonzero in the polynomial integral domain. Exactly four factors
meet each site; thus \(B_{L,N}\) has degree four there and is a
vector of \(\mathcal H_\Lambda\).
It is divisible by \(q_e\) for every \(e\), so
Lemma~\ref{l:edge} gives \(P_eB_{L,N}=0\).

For any \(\varphi\), orthogonality of the projections gives
\[
 \langle\varphi,H_{L,N}\varphi\rangle
       =\sum_e\|P_e\varphi\|^2.
\]
Consequently \(\ker H_{L,N}=\bigcap_e\ker P_e\).
Apply Lemma~\ref{l:divisibility} to a vector in this intersection
with \(R=\Lambda\). The quotient by \(B_{L,N}\) is independent
of every site variable and is a scalar. This proves both inclusions
in the kernel equality. Reversing any chosen orientations multiplies
\(B_{L,N}\) by \((-1)^k\), where \(k\) is the number reversed,
and therefore preserves its span.
\end{proof}

\begin{lemma}\label{l:star}
Let \(h_\star=\sum_{j=1}^4P_{\{0,j\}}\) on \(V_4^{\otimes5}\).
Its smallest positive eigenvalue \(g_\star\) exists and is positive.
For \(L,N\ge3\), set \(H_x=\sum_{e\ni x}P_e\) and
\(\mathcal S_x=\ker H_x\). Then
\[
 H_x\ge g_\star(I-\Pi_{\mathcal S_x}),\qquad
 \sum_xH_x=2H_{L,N}.
\]
\end{lemma}

\begin{proof}
The operator \(h_\star\) is positive semidefinite on the
\(5^5\)-dimensional Hilbert space. It is nonzero: the tensor
\((z_2^4)^{\otimes5}\) has each endpoint pair in the
one-dimensional weight-four subspace used in Lemma~\ref{l:edge},
so every one of the four projections fixes it. Its eigenvalue for
\(h_\star\) is four. The finite set of positive eigenvalues is
therefore nonempty and has a strictly positive minimum \(g_\star\).
The spectral theorem gives
\(h_\star\ge g_\star(I-\Pi_{\ker h_\star})\).

The center \(x\) and its four neighbors are distinct, so relabelling
those five tensor factors identifies \(H_x\) with
\(h_\star\) tensored with the identity on the other sites.
Edges between the neighbors do not occur in \(H_x\) and do not
affect this identification. Tensoring the preceding inequality
proves the first assertion for \(H_x\). Each unoriented edge has
two different endpoints and occurs once at each endpoint in
\(\sum_xH_x\), proving the second assertion.
\end{proof}

\subsection{Finite frames and their separation}

The symbols \(D\), \(\phi\), \(r\), and \(\mathcal M\) have the
meaning fixed in Definition~\ref{d:moment}. In particular,
\(\phi=(1+\sqrt5)/2\), \(r^2=\phi+2\), and
\(\mathcal M\) is the uniform average on the twelve displayed
icosahedral vertices. We use its degree-four moment identities and
polynomial covariance from Lemma~\ref{l:moments}.

\begin{definition}\label{d:frame}
A frame is a finite set \(A\subset\mathbb S^2\) with positive weights
\(w(a)\) summing to one such that
\(\sum_{a\in A}w(a)p(a)=\mathcal M(p)\)
for every real polynomial \(p\) of degree at most four.
\end{definition}

The defining identity also holds for complex polynomials by applying
it to their real and imaginary parts.

\begin{lemma}\label{l:tight}
For every frame \((A,w)\) and every \(p\in V_4\),
\[
 \sum_{a\in A}5w(a)|p(\zeta(a))|^2=\|p\|_{V_4}^2.
\]
For site frames \((D_x,w_x)\), set \(\Omega=\prod_xD_x\),
\(\pi(s)=\prod_x5w_x(s_x)\), and
\(\iota\varphi(s)=\varphi((\zeta(s_x))_x)\).
Then \(\iota:\mathcal H_\Lambda\to\ell^2(\Omega,\pi)\) is an isometry.
\end{lemma}

\begin{proof}
Let \(m_j(z)=z_1^{4-j}z_2^j\), \(0\le j\le4\).
For \(j\le k\), the spinor projection identity gives
\[
 \overline{m_j(\zeta(a))}\,m_k(\zeta(a))
 =P_{jk}(a):=
 \frac{(1+a_3)^{4-k}(1-a_3)^j(a_1+i a_2)^{k-j}}{16}.
\]
Indeed, the paired factors are
\(|\zeta_1|^2=(1+a_3)/2\),
\(|\zeta_2|^2=(1-a_3)/2\), and
\(\overline{\zeta_1}\zeta_2=(a_1+i a_2)/2\).
For \(j>k\), use \(P_{jk}=\overline{P_{kj}}\).
Every \(P_{jk}\) has degree at most four. If \(j<k\), rotation
through \(\pi/(k-j)\) about the third coordinate multiplies
\(P_{jk}\) by \(-1\). Polynomial covariance in
Lemma~\ref{l:moments} implies \(\mathcal M(P_{jk})=0\).
The same follows by conjugation for \(j>k\).

For diagonal entries only the moments of \(t=a_3\) are needed:
\(\mathcal M(1)=1\), \(\mathcal M(t)=\mathcal M(t^3)=0\),
\(\mathcal M(t^2)=1/3\), and \(\mathcal M(t^4)=1/5\).
Expanding the five polynomials yields
\[
 \begin{array}{c|ccccc}
 j&0&1&2&3&4\\ \hline
 \mathcal M(P_{jj})&1/5&1/20&1/30&1/20&1/5
 \end{array}
 \qquad
 5\mathcal M(P_{jj})=\binom4j^{-1}.
\]
For example, the three different even parts of their numerators
are \(1+6t^2+t^4\), \(1-t^4\), and
\(1-2t^2+t^4\), respectively. The frame identity for these
degree-four polynomials now gives
\[
 \sum_{a\in A}5w(a)
       \overline{m_j(\zeta(a))}m_k(\zeta(a))
       =\begin{cases}\binom4j^{-1},&j=k,\\0,&j\ne k.\end{cases}
\]
Expanding \(p\) in the monomials proves its norm identity with
exactly the normalization of Definition~\ref{d:poly}.

For a finite collection of sites, the evaluation map is the tensor
product of the single-site evaluation maps. The product monomials
form an orthogonal basis, and the finite sums of their pairwise
inner products factor into the single-site sums just computed.
Their tensor inner products are therefore preserved, which proves
the isometry on all of \(\mathcal H_\Lambda\).
\end{proof}

The weight \(\pi\) in this lemma is a positive reference measure
of total mass \(5^{|\Lambda|}\), rather than a probability.
The use of rotations in the proof was solely on polynomials of
degree at most four.

\begin{lemma}\label{l:threeframes}
Let
\[
 T=\begin{pmatrix}99/101&-20/101&0\\20/101&99/101&0\\0&0&1\end{pmatrix},
 \quad D^{(j)}=T^jD\ (j=0,1,2),\quad
 b=\frac{5-\sqrt5}{505},\quad r_b=\frac b{2-b}.
\]
The matrix \(T\) is in \(\mathrm{SO}(3)\).
Each \(D^{(j)}\), with uniform weights \(1/12\), is an antipodal frame.
For \(i\ne j\), \(a\in D^{(i)}\), \(a'\in D^{(j)}\),
\[
 0<b<1,\qquad b\le1\pm a\cdot a'\le2-b,\qquad0<r_b<1.
\]
The common lower bound \(b\) is attained, and distinct frames are
disjoint even up to antipodal signs.
\end{lemma}

\begin{proof}
The equality \(99^2+20^2=101^2\) proves \(T^TT=I\) and
\(\det T=1\). The set \(D\) consists of twelve distinct unit
vectors and is antipodal by Lemma~\ref{l:moments}.
For a polynomial \(p\) of degree at most four,
\[
 \frac1{12}\sum_{a\in D^{(j)}}p(a)
       =\mathcal M(p\circ T^j)=\mathcal M(p).
\]
Here \(T^j\) is bijective and covariance applies on precisely
the stated degree range. Thus each rotated set is an antipodal
frame of twelve letters.

Write \(\vartheta\) for the rotation angle of \(T\), so that
\(\cos\vartheta=99/101\) and \(\sin\vartheta=20/101\).
We have
\[
 1-\cos\vartheta=\frac2{101},\qquad
 1-\cos(2\vartheta)=\frac{800}{10201}\ge\frac2{101},
 \qquad
 b=\frac{2}{101(\phi+2)}.
\]
The last identity follows by rationalizing
\(\phi+2=(5+\sqrt5)/2\). In particular \(0<b<1\).
By applying a common inverse rotation and, if needed, swapping
the two vectors, it suffices to bound
\(1-u\cdot T^k v\) for \(u,v\in D\) and \(k\in\{1,2\}\).

First let \(u=v\). Direct multiplication of the planar rotation
matrix gives
\[
 1-u\cdot T^ku
       =(1-\cos(k\vartheta))(u_1^2+u_2^2).
\]
The three possible values of \(u_1^2+u_2^2\), from the three
coordinate families of \(D\), are
\(1/r^2,1,\phi^2/r^2\). Since \(\phi>1\), their minimum is
\(1/r^2\). It follows that \(1-u\cdot T^ku\ge b\).
Equality holds for \(k=1\) and \(u=(0,1,\phi)/r\).

Next suppose \(u\ne v\). For vertices in the same coordinate
family, their scalar product has the form
\((\epsilon\epsilon'+\phi^2\delta\delta')/r^2\) with
\(\epsilon,\epsilon',\delta,\delta'\in\{\pm1\}\).
Apart from the equal-vertex case, its values are
\(-1\) and \(\pm\phi/r^2\). For vertices in different
families, the only shared nonzero coordinate has product
\(\pm\phi\), so the scalar product is \(\pm\phi/r^2\).
The identity \(\phi/r^2=1/\sqrt5\) thus proves
\[
 u\cdot v\in\{-1,1/\sqrt5,-1/\sqrt5\},\qquad
 |u-v|^2\ge2-\frac2{\sqrt5}>1.
\]
On the first two coordinates the operator norm of a rotation minus
the identity is \(\sqrt{2(1-\cos(k\vartheta))}\), while it is
zero on the third coordinate. Hence
\[
 \|T-I\|=\frac2{\sqrt{101}}\le\frac{40}{101},\qquad
 \|T^2-I\|=\frac{40}{101}<\frac25.
\]
The triangle inequality gives \(|u-T^kv|>3/5\), and therefore
\[
 1-u\cdot T^kv=\tfrac12|u-T^kv|^2>\frac9{50}>b.
\]
For the last inequality, \(\phi+2>2\) implies
\(b<1/101<9/50\). This proves \(1-a\cdot a'\ge b\)
in every case. Replace \(a'\) by \(-a'\), which belongs to the
same rotated frame, to obtain \(1+a\cdot a'\ge b\) as well.
These two lower bounds give both upper bounds \(2-b\).
The equality case already displayed attains \(b\); changing the
sign of its second vector also attains the bound with the other
choice of sign. An equality \(a=\pm a'\) would make one of
\(1\pm a\cdot a'\) zero, contradicting \(b>0\).
Finally, \(0<b<1\) implies \(0<b/(2-b)<1\).
\end{proof}

\begin{lemma}\label{l:colour}
For \(m\ge3\) and representatives \(0\le i<m\), put
\[
 a_m(i)=\begin{cases}2,&m\text{ odd and }i=m-1,\\i\bmod2,&\text{otherwise}.
 \end{cases}
\]
Then \(\kappa(i,j)=a_L(i)+a_N(j)\pmod3\) is a proper
three-coloring of \(\Lambda_{L,N}\) for every \(L,N\ge3\).
\end{lemma}

\begin{proof}
If \(m\) is even, the two colors zero and one alternate along
the cycle, and the wraparound edge joins the odd index \(m-1\)
to index zero. If \(m\) is odd, they alternate on indices
\(0,\ldots,m-2\); the last of these indices is odd and has
color one. The remaining vertex \(m-1\) has color two and is
adjacent to that vertex and to index zero, whose color is zero.
This covers \(m=3\) as well. Thus every cycle edge changes
\(a_m\) by a nonzero element of \(\mathbb Z_3\).
A torus edge changes exactly one summand in the formula for
\(\kappa\), leaving the other fixed, so its endpoint colors
are different modulo three.
\end{proof}

\begin{lemma}\label{l:qmod}
For unit spinors \(\zeta,\zeta'\) with
\(\zeta\zeta^*=(I+a\cdot\sigma)/2\) and
\(\zeta'\zeta'^*=(I+a'\cdot\sigma)/2\),
\[
 |\zeta_1\zeta'_2-\zeta_2\zeta'_1|^2=(1-a\cdot a')/2 .
\]
\end{lemma}

\begin{proof}
Expansion of the two-by-two determinant gives the Gram identity
\[
 |\zeta_1\zeta'_2-\zeta_2\zeta'_1|^2
       =\|\zeta\|^2\|\zeta'\|^2-|\zeta^*\zeta'|^2
       =1-|\zeta^*\zeta'|^2.
\]
The displayed Pauli matrices satisfy
\(\operatorname{tr}\sigma_j=0\) and
\(\operatorname{tr}(\sigma_j\sigma_k)=2\delta_{jk}\).
Consequently
\[
 |\zeta^*\zeta'|^2
 =\operatorname{tr}(\zeta\zeta^*\zeta'\zeta'^*)
 =\frac14\operatorname{tr}\bigl((I+a\cdot\sigma)
                               (I+a'\cdot\sigma)\bigr)
 =\frac{1+a\cdot a'}2.
\]
Substitution proves the assertion.
\end{proof}

\subsection{Frame specifications and the quantum comparison}

\begin{definition}\label{d:framemodel}
Given any proper three-coloring \(\kappa\) of a plane or torus,
put \(D_x=D^{(\kappa(x))}\), with weights \(1/12\).
On the torus define
\[
 \mu^{\mathrm{af}}_{\kappa,L,N}(s)=
 \frac1{Z^{\mathrm{af}}}\prod_x\frac1{12}
       \prod_{\{x,y\}\in E_{L,N}}(1-s_x\cdot s_y).
\]
On the plane, for finite \(W\subset\mathbb Z^2\) and
\(\eta_z\in D_z\) at exterior neighbor sites, define
\[
 \mu^{+,\eta}_{\kappa,W}(n)=\frac1{Z^{+,\eta}_{\kappa,W}}
 \prod_{x\in W}\frac1{12}
 \prod_{\substack{\{x,y\}\subset W\\d(x,y)=1}}(1+n_x\cdot n_y)
 \prod_{\substack{x\in W,z\notin W\\d(x,z)=1}}(1+n_x\cdot\eta_z).
\]
Only the finitely many needed exterior letters are specified.
These laws have full support. Write \(E_U\) for conditioning off \(U\),
\(\operatorname{Var}_U f=E_U|f-E_Uf|^2\), and
\[
 \mathcal D_\mu(f)=\sum_x\mu(\operatorname{Var}_x f),\qquad
 \lambda(\mu)=\inf_{\operatorname{Var}_\mu f>0}
             \frac{\mathcal D_\mu(f)}{\operatorname{Var}_\mu f}.
\]
All generators have one update per physical site. A gap lower bound is
always equivalently a Poincare inequality, including one-point spaces.
\end{definition}

Indeed, every interaction factor in this definition lies in
\([b,2-b]\), by proper coloring and Lemma~\ref{l:threeframes}.
There are finitely many configurations, each of strictly positive
weight, so every normalizer is finite and positive. For
\(W=\varnothing\) the normalizer is one and the configuration
space has one point.
On a positive finite law the conditional expectation \(E_U\) is
the orthogonal projection onto the functions of the letters off
\(U\): this follows by summing on each fixed exterior fiber.
In particular,
\[
 \mu(\operatorname{Var}_U f)=\|(I-E_U)f\|_{L^2(\mu)}^2,
 \qquad \mathcal L_\mu=\sum_x(E_x-I).
\]
If \(\mathcal D_\mu(f)=0\), full support makes \(f\) constant
on every single-site fiber. Any two product configurations can be
joined by changing their coordinates one at a time, so \(f\) is
constant. The finite-dimensional spectral theorem consequently
gives \(\lambda(\mu)>0\) whenever the configuration space has
at least two points. On a one-point space every Poincare inequality
holds by the convention of Section~1.

We abbreviate the torus law by \(\mu^{\rm af}\) and the valence-bond
state by \(B\). Every evaluation of an incident determinant is
nonzero: its squared modulus is at least \(b/2\) by
Lemmas~\ref{l:threeframes} and~\ref{l:qmod}.

\begin{lemma}\label{l:quantumfibre}
Fix \(x\), write \(q_e=\epsilon_{e,x}q_{xy}\) for \(e=\{x,y\}\),
and set
\[
 c_x=\prod_{e\ni x}\epsilon_{e,x},\quad
 \beta_x(t;\sigma)=\prod_{y\sim x}q_{xy}(\zeta(t),\zeta(\sigma_y)),
 \quad \rho_x=c_x\prod_{e\not\ni x}q_e .
\]
All products are evaluated at the corresponding frame spinors.
Then \(\iota B=\beta_x\rho_x\), and if
\(\psi=(\prod_{e\ni x}q_e)\chi\) with \(\chi\) independent of \(z_x\),
then \(\iota\psi=c_x\beta_x\,\iota\chi\).
Let \(\mathcal P_x\) be the fiberwise orthogonal projection in
\(\ell^2(\Omega,\pi)\) onto the span of \(\beta_x\), and put
\(Uf=f\,\iota B/\|B\|\).
The map \(U:L^2(\mu^{\rm af})\to\ell^2(\Omega,\pi)\) is unitary,
\(UE_x=\mathcal P_xU\), and
\[
 \|(I-\mathcal P_x)\iota\varphi\|
            \le\|(I-\Pi_{\mathcal S_x})\varphi\|.
\]
\end{lemma}

\begin{proof}
Here \(\epsilon_{e,x}=1\) when the chosen orientation of \(e\)
starts at \(x\), and \(\epsilon_{e,x}=-1\) otherwise, so
\(c_x\in\{1,-1\}\). Separating the four incident edge factors
in the product for \(B\) gives \(\iota B=\beta_x\rho_x\).
Evaluation of a product is the product of its evaluations, which
also gives \(\iota\psi=c_x\beta_x\,\iota\chi\) in the stated
case. The notation \(\iota\chi\) in this equality means evaluation
of the residual polynomial; no assertion about an isometry for its
different degrees is used. Both \(\rho_x\) and this evaluation
are independent of the coordinate at \(x\).

We first verify the probability identity that is needed for
unitarity. Lemma~\ref{l:qmod} and the product reference weights give
\begin{equation}\label{e:frame-ground-weight}
 \pi(s)|\iota B(s)|^2
 =5^{|\Lambda|}2^{-|E_{L,N}|}
       \prod_x\frac1{12}
       \prod_{\{u,v\}\in E_{L,N}}(1-s_u\cdot s_v).
\end{equation}
By Lemma~\ref{l:tight}, summing the left side gives \(\|B\|^2\).
Thus
\(\|B\|^2=5^{|\Lambda|}2^{-|E_{L,N}|}Z^{\rm af}\), and
\begin{equation}\label{e:frame-ground-law}
 \mu^{\rm af}(s)=\frac{\pi(s)|\iota B(s)|^2}{\|B\|^2}.
\end{equation}
The product \(\iota B\) is nowhere zero. Equation
\eqref{e:frame-ground-law} proves that \(U\) preserves the norm,
and its inverse on the full finite function space is
\(h\mapsto\|B\|h/(\iota B)\). Therefore it is unitary.

Fix an exterior configuration \(s_{\Lambda\setminus\{x\}}\)
and let \(\sigma\) denote its neighbor letters. The conditional
probabilities, by either the factorization or
\eqref{e:frame-ground-law}, are
\[
 \mu^{\rm af}(t\mid s_{\Lambda\setminus\{x\}})
 =\frac{5w_x(t)|\beta_x(t;\sigma)|^2}
        {\sum_{u\in D_x}5w_x(u)|\beta_x(u;\sigma)|^2}.
\]
The denominator is strictly positive. On the same fiber the
orthogonal projection formula in \(\ell^2(D_x,5w_x)\) is
\[
 (\mathcal P_xh)(t,s_{\Lambda\setminus\{x\}})
 =\beta_x(t;\sigma)
  \frac{\sum_{u\in D_x}5w_x(u)\overline{\beta_x(u;\sigma)}
                         h(u,s_{\Lambda\setminus\{x\}})}
       {\sum_{u\in D_x}5w_x(u)|\beta_x(u;\sigma)|^2}.
\]
Since the outside factor of \(\pi\) is constant on that fiber,
these orthogonal fiber projections combine to an orthogonal
projection on \(\ell^2(\Omega,\pi)\).
Substitute \(h=Uf=\beta_x\rho_x f/\|B\|\) in the formula.
The result is \(\beta_x\rho_x(E_xf)/\|B\|=U(E_xf)\), proving
the intertwining.

If \(\psi\in\mathcal S_x\), positivity of its four projection
terms implies \(P_e\psi=0\) for all \(e\ni x\).
Lemma~\ref{l:divisibility}, with \(R=\{x\}\), gives the
factorization of \(\psi\) used in the statement with a quotient
independent of \(z_x\). Its evaluated fiber is consequently a
multiple of \(\beta_x\). This proves the inclusion
\(\iota\mathcal S_x\subset\operatorname{Ran}\mathcal P_x\).
It follows that
\[
 \begin{split}
 \|(I-\mathcal P_x)\iota\varphi\|
 &=\|(I-\mathcal P_x)\iota(I-\Pi_{\mathcal S_x})\varphi\|\\
 &\le\|\iota(I-\Pi_{\mathcal S_x})\varphi\|
  =\|(I-\Pi_{\mathcal S_x})\varphi\|.
 \end{split}
\]
The inequality is the contraction property of an orthogonal
projection, and the last equality is the physical-space isometry
of Lemma~\ref{l:tight}.
\end{proof}

\begin{theorem}\label{t:QD}
For every \(L,N\ge3\) and every proper torus three-coloring,
\[
 \mu^{\mathrm{af}}(s)=
 \frac{\pi(s)|\iota B_{L,N}(s)|^2}{\|B_{L,N}\|^2},\qquad
 H_{L,N}\big|_{B_{L,N}^{\perp}}\ge
       \frac{g_\star}{2}\lambda(\mu^{\mathrm{af}})I .
\]
In particular
\(\operatorname{gap}(H_{L,N})\ge(g_\star/2)\lambda(\mu^{\mathrm{af}})\).
\end{theorem}

\begin{proof}
The probability identity was proved in
\eqref{e:frame-ground-law}. For \(\varphi\perp B\), set
\(f=U^{-1}\iota\varphi\). Since
\(U1=\iota B/\|B\|\) and both maps preserve inner products,
\[
 \mu^{\rm af}(f)
 =\langle U1,\iota\varphi\rangle
 =\frac{\langle B,\varphi\rangle}{\|B\|}=0,\qquad
 \operatorname{Var}_{\mu^{\rm af}}f=\|\varphi\|^2.
\]
Using the Poincare definition of \(\lambda\), the intertwining
and norm inequality of Lemma~\ref{l:quantumfibre}, and
Lemma~\ref{l:star}, we obtain
\[
 \begin{split}
 \lambda(\mu^{\rm af})\|\varphi\|^2
 &\le\sum_x\|(I-E_x)f\|_{L^2(\mu^{\rm af})}^2\\
 &=\sum_x\|(I-\mathcal P_x)\iota\varphi\|_{\ell^2(\pi)}^2\\
 &\le\sum_x\langle\varphi,(I-\Pi_{\mathcal S_x})\varphi\rangle\\
 &\le\frac1{g_\star}\sum_x\langle\varphi,H_x\varphi\rangle
  =\frac2{g_\star}\langle\varphi,H_{L,N}\varphi\rangle.
 \end{split}
\]
For \(\varphi=0\) the same inequality is immediate.
The space \(B^\perp\) is invariant under \(H_{L,N}\), since
\(H_{L,N}\) is self-adjoint and kills \(B\).
Proposition~\ref{p:toruskernel} identifies its spectrum with the
positive spectrum of \(H_{L,N}\), proving the last assertion.
\end{proof}

\subsection{Finite energy and the gauge on a lifted square}

\begin{lemma}\label{l:planeenergy}
All the frame specifications above are consistent under conditioning
and have nearest-neighbor Markov dependence. Uniformly over proper
colorings, finite volumes, admissible frame-valued boundaries, and
fully specified one-site fibers, every letter has conditional
probability at least
\[
 \alpha=\frac1{12}r_b^4\in(0,1].
\]
The same lower bound holds after conditioning on fewer letters.
Any prescribed assignment on \(k\) sites, under arbitrary consistent
partial conditioning off those sites, has conditional probability
at least \(\alpha^k\).
For a finite plane block \(Q\), the single-site heat bath of every conditional block law in \(Q\) has Poincare gap at least
\(r_b^{|E(Q)|}\), where \(E(Q)\) is its internal edge set.
\end{lemma}

\begin{proof}
Let \(U\) be a subset of the unfixed sites in a finite specification
and fix the letters on its complement. Every site or edge factor
not meeting \(U\) is then constant and cancels between the
conditional weight and its normalizer. The remaining factors are
exactly the site factors on \(U\), the internal edge factors of
\(U\), and the crossing factors with the newly fixed neighbor
letters. They have the same sign of interaction as the initial
specification. This proves consistency, both for the torus law and
for the plane laws, and shows that only exterior nearest-neighbor
letters affect the conditional law. The assertion also holds for
\(U=\varnothing\), whose conditional space has one point.

For a fully specified fiber at a site \(x\), the unnormalized
weight of a proposed letter \(t\in D_x\) is \(1/12\) times
at most four factors, each in \([b,2-b]\). If their number is
\(d_x\le4\), the numerator is at least \(b^{d_x}/12\),
and the sum of all twelve weights is at most \((2-b)^{d_x}\).
Thus
\[
 \Pr(s_x=t\mid\hbox{all other unfixed letters})
 \ge\frac1{12}r_b^{d_x}\ge\frac1{12}r_b^4=\alpha.
\]
Lemma~\ref{l:threeframes} gives \(0<r_b<1\), hence
\(0<\alpha\le1\). Conditioning on fewer letters takes a finite
weighted average of the fully specified conditional probabilities;
every summand has this same lower bound.

More explicitly, let \(A=\{x_1,\ldots,x_k\}\) consist of
distinct sites and prescribe \(t_i\in D_{x_i}\). Condition on
any positive event specifying letters off \(A\). At stage \(i\)
of the chain rule, that event and the previous \(i-1\) prescribed
letters condition only off \(x_i\), so the conditional probability
of \(t_i\) is at least \(\alpha\). The preceding lower bounds
ensure that every intermediate event has positive probability.
Multiplying gives \(\alpha^k\). For \(k=0\), the probability
and the asserted bound are both one. The argument also applies to
a positive event determined by fewer outside letters, by averaging.

For the block-gap assertion, fix a finite plane block \(Q\) and
all its exterior neighbor letters \(\eta\). Absorb every fixed
boundary factor into its endpoint site weight by setting
\[
 h_x(t)=\frac1{12}
       \prod_{\substack{z\notin Q\\d(x,z)=1}}
                    (1+t\cdot\eta_z),\qquad
 \nu_x(t)=\frac{h_x(t)}{\sum_{u\in D_x}h_x(u)},\qquad
 \nu_Q=\bigotimes_{x\in Q}\nu_x.
\]
Every \(h_x(t)\) is positive, since all its boundary factors are
in \([b,2-b]\). Thus \(\nu_Q\) is a strictly positive product
probability on the finite product \(\prod_{x\in Q}D_x\). Let
\[
 R(n)=\prod_{\{x,y\}\in E(Q)}(1+n_x\cdot n_y),\qquad
 e=|E(Q)|.
\]
The conditional block law is \(d\mu=(R/\nu_Q(R))\,d\nu_Q\),
and \(b^e\le R\le(2-b)^e\). In particular,
\(0<\nu_Q(R)<\infty\), and its density has the bounds
\[
 0<m_Q:=\frac{b^e}{\nu_Q(R)}
 \ \le\ \frac{R(n)}{\nu_Q(R)}
 \ \le\ M_Q:=\frac{(2-b)^e}{\nu_Q(R)}<\infty.
\]
Apply the density comparison of Lemma~\ref{l:density}, proved in
Appendix~A, with its product reference law \(\pi=\nu_Q\),
density \(\varrho=R/\nu_Q(R)\), and constants \(m=m_Q\),
\(M=M_Q\). All its hypotheses have just been verified. Its
conditional-variance comparison and product Poincare inequality give,
for every complex \(f\),
\[
 \mathcal D_\mu(f)
 \ge m_Q\mathcal D_{\nu_Q}(f)
 \ge m_Q\operatorname{Var}_{\nu_Q}f
 \ge\frac{m_Q}{M_Q}\operatorname{Var}_\mu f.
\]
Consequently the Poincare gap is at least
\(m_Q/M_Q=r_b^e\). Only internal edges occur in this ratio;
all fixed boundary factors have been retained in \(\nu_Q\).
If some sites of \(Q\) are already fixed, apply this construction
to the free subset \(U\subset Q\). Since
\(E(U)\subset E(Q)\) and \(0<r_b<1\), the resulting bound
\(r_b^{|E(U)|}\) is at least the displayed bound for \(Q\).
If the free subset is empty, the Poincare inequality holds because
its variance is zero. In particular the argument includes empty
blocks and blocks with no internal edges.
\end{proof}

The finite-energy bound in this lemma concerns admissible frame
letters on neighbor sites. Arbitrary unit-vector pins, which are
used in the separate response estimates, are not included in its
quantifiers.

\begin{lemma}\label{l:boxlift}
Let \(Q'=a+\{-M,\ldots,M\}^2\subset\mathbb Z^2\), \(M\ge1\),
and \(L,N\ge2M+3\). Reduction modulo the periods is injective
on the \(\ell^\infty\) one-neighborhood of \(Q'\), preserves exactly
the internal nearest-neighbor edges of \(Q'\), and identifies all its
crossing bonds with the torus crossing bonds of its image \(Q\).
For every torus outside configuration \(\sigma\), pull back the
coloring and set
\[
 n_z=(-1)^{z_1+z_2}s_{p(z)}\quad(z\in Q'),\qquad
 \eta_z=(-1)^{z_1+z_2}\sigma_{p(z)}
       \quad(z\notin Q',\ d(z,Q')=1).
\]
This coordinatewise bijection carries the conditional
antiferromagnetic law in \(Q\) to
\(\mu^{+,\eta}_{\kappa\circ p,Q'}\), intertwines every site update,
and preserves Hamming distances between configurations in each fiber.
\end{lemma}

\begin{proof}
Let \(p:\mathbb Z^2\to\mathbb Z_L\times\mathbb Z_N\) be
coordinatewise reduction. The indicated one-neighborhood is
\[
 A'=a+\{-M-1,\ldots,M+1\}^2.
\]
The difference between two of its first coordinates has absolute
value at most \(2M+2<L\), and the analogous bound is
\(2M+2<N\) for second coordinates. If their reductions agree,
both differences are multiples of their periods, so both are zero.
This proves injectivity on \(A'\).

Plane nearest neighbors reduce to distinct torus nearest neighbors.
Conversely, if \(z,w\in Q'\) have nearest-neighbor reductions,
the congruent coordinate differences are \(\pm1\) and zero,
in one order. A difference in \([-2M,2M]\) congruent to zero
modulo its period must be zero. A difference in that interval
congruent to \(1\) or \(-1\) must be that very integer:
every other representative has absolute value at least
\(\min\{L,N\}-1\ge2M+2\). Thus \(d(z,w)=1\), which
proves the exact correspondence of internal edges.

For a torus edge leaving \(Q=p(Q')\), lift its endpoint in \(Q\)
uniquely to \(z\in Q'\). Its other endpoint has the form
\(p(z+v)\), with a unique
\(v\in\{(1,0),(-1,0),(0,1),(0,-1)\}\); the four torus
neighbors are distinct. Then \(z+v\in A'\), and it is outside
\(Q'\) since the torus endpoint is outside \(Q\).
Conversely, a plane crossing bond from \(z\in Q'\) to
\(w\notin Q'\) has \(w\in A'\); injectivity on \(A'\)
ensures \(p(w)\notin Q\). These operations are inverse on
crossing bonds. They also show that all required exterior sites
are represented without identifications.

The function \(\kappa\circ p\) is a proper plane coloring:
each plane nearest-neighbor edge maps to a torus edge whose
colors differ. Put \(\varepsilon_z=(-1)^{z_1+z_2}\) for the
plane sites involved in the lift. Each alphabet is antipodal with
uniform weights, so multiplication by \(\varepsilon_z\) is a
weight-preserving bijection on its alphabet. The displayed \(\eta\)
are consequently admissible exterior frame letters.
For a plane internal or crossing bond \(\{z,w\}\),
\(\varepsilon_z\varepsilon_w=-1\). Substituting
\(s_{p(z)}=\varepsilon_z n_z\) at a free site and
\(\sigma_{p(w)}=\varepsilon_w\eta_w\) at a fixed site changes
its torus factor as follows:
\[
 1-s_{p(z)}\cdot s_{p(w)}=1+n_z\cdot n_w,
 \qquad
 1-s_{p(z)}\cdot\sigma_{p(w)}=1+n_z\cdot\eta_w.
\]
Factors wholly outside \(Q\) cancel in the conditional torus
law. The exact internal and crossing-bond correspondences show
that all remaining factors are precisely those in
\(\mu^{+,\eta}_{\kappa\circ p,Q'}\). Site weights are unchanged,
and the configuration bijection makes their normalizing sums
equal. This proves the asserted transport of laws.

For \(x=p(z)\), let \(Jf(s)=f((\varepsilon_w s_{p(w)})_{w\in Q'})\)
be pullback on the corresponding conditional function spaces.
It is unitary by the transported law. In a single-site fiber all
other coordinates are fixed, and summing over the letter at \(x\)
is exactly summing over its image letter at \(z\) with the
transported conditional probabilities. Therefore
\(E_x^{\rm af}J=J E_z^+\). This proves the update intertwining.
Finally \(s_{p(z)}=t_{p(z)}\) if and only if
\(\varepsilon_zs_{p(z)}=\varepsilon_zt_{p(z)}\), so the
coordinate bijection preserves the number of unequal coordinates.
All signs used here are defined on plane lifts; no sign function
on the torus is needed.
\end{proof}
